%======================================================================
\chapter{Preliminaries}
\label{ch:prelim}
%======================================================================

This chapter assumes no background beyond ordinary programming. Every cryptographic idea the
protocol relies on is introduced here with a concrete example, so that Chapters~3 onwards
read without interruption.

\section{One-way hash functions}

A cryptographic hash function maps an input of \emph{any} length to a fixed-length
fingerprint:
\begin{equation}
H : \{0,1\}^{*} \longrightarrow \{0,1\}^{n}
\end{equation}
SHA-256 produces $n = 256$ bits, conventionally written as 64 hexadecimal characters.

\begin{lstlisting}[style=plain,caption={Avalanche: a one-character change alters roughly half the output bits.}]
SHA256("underwater")  ->  8e4c9d6f2b1a7053c8e4...   (64 hex characters)
SHA256("underwatex")  ->  41f7a02c9db35e18770c...   (no visible relation)
\end{lstlisting}

Three properties matter for this protocol:

\begin{itemize}
\item \textbf{Pre-image resistance (one-wayness).} Given the output, recovering the input is
      infeasible; exhaustive search costs about $2^{256}$ attempts.
\item \textbf{Collision resistance.} Two distinct inputs with the same output cannot be
      found in practice.
\item \textbf{Avalanche.} Flipping one input bit flips about half the output bits, so the
      digest carries no partial information about the input.
\end{itemize}

The protocol uses hashing for two distinct jobs: turning a public key into a short
unforgeable \emph{identity}, and condensing a raw ECDH shared point into a usable
\emph{symmetric key}.

\section{Symmetric versus asymmetric encryption}

\textbf{Symmetric} encryption uses the same key to encrypt and decrypt. AES is the standard.
It is fast and cheap in both time and energy --- exactly what a battery-powered sensor
needs. Its difficulty is key distribution: both parties must already share the key, over a
channel an adversary is listening to.

\textbf{Asymmetric} (public-key) encryption uses a key \emph{pair}: a public key published
freely, and a private key never revealed. Anything encrypted under the public key can be
opened only with the matching private key. This solves distribution but costs two to three
orders of magnitude more computation.

Practical systems, including this one, are \textbf{hybrid}: asymmetric mathematics
establishes a shared secret, and fast symmetric encryption then protects the actual
messages.

\section{Elliptic curve cryptography}

An elliptic curve over a finite field $\mathbb{F}_p$ is the set of points $(x,y)$ satisfying
\begin{equation}
y^{2} \equiv x^{3} + ax + b \pmod{p}
\end{equation}
together with a point at infinity. Our implementation uses \textbf{secp256r1}
(NIST P-256), for which $p \approx 2^{256}$.

Curve points can be added by a geometric rule, and adding a base point $G$ to itself $k$
times is written $k \cdot G$ --- \textbf{scalar multiplication}. This single operation is the
entire engine of ECC, because it is easy in one direction and infeasible in the other:

\begin{itemize}
\item Computing $k \cdot G$ from $k$ takes about 256 double-and-add steps --- microseconds.
\item Recovering $k$ from $k \cdot G$ is the \textbf{Elliptic Curve Discrete Logarithm
      Problem} (ECDLP). The best known attack needs roughly $2^{128}$ operations on a
      256-bit curve.
\end{itemize}

A \textbf{private key} is therefore just a random integer $k \in [1, n-1]$, and the matching
\textbf{public key} is the curve point $k \cdot G$. Publishing that point discloses nothing
about $k$.

\begin{keybox}
\textbf{Why ECC and not RSA here.} ECC reaches 128-bit security with a 256-bit key; RSA
needs 3072 bits for the same. On a 10\,kbps acoustic link, transmitting a 3072-bit key costs
307\,ms against 26\,ms for ECC --- and that gap is paid on every hop, every session, out of
a battery that cannot be replaced.
\end{keybox}

\subsection{The paper's pentatope extension (PEECC)}

The base paper proposes an extension it calls \emph{pentatope elliptic curve cryptography}.
Instead of a two-dimensional point group, it defines a five-dimensional group
$G_5 = \{(x_1,\dots,x_5) \mid x_i \in \mathbb{F}_p\}$ over the curve
$\psi = X^3+Y^3+Z^3+W^3+T^3 = d\,XYZWT$. Key generation follows standard ECC but folds in
the \emph{pentatope number}, the fifth entry of a Pascal's-triangle diagonal,
\begin{equation}
S = \frac{S(S+1)(S+2)(S+3)}{24},
\qquad K = S^{2} + S + 41,
\end{equation}
which the authors use to derive a key value internally rather than importing an externally
computed shared secret. Their claim is that the discrete logarithm problem in $G_5$ has
solving complexity around $2^{5n/2}$ against $2^{2n/2}$ for a conventional curve, at a
measured $1.3\times$ scalar-multiplication overhead on an ARM Cortex-A53, and that the
number of global parameters falls from six to five.

\begin{warnbox}
\textbf{Scope note.} Our implementation uses standard secp256r1, \emph{not} PEECC. The
pentatope construction is not accompanied in the paper by a reference implementation or
concrete curve parameters, so reproducing it was outside what we could verify. Everything we
measure is therefore a lower bound on the paper's claimed security, not a reproduction of
it. This is discussed in Section~\ref{sec:deviations}.
\end{warnbox}

\section{ECDH: agreeing on a secret in public}

Elliptic Curve Diffie--Hellman lets two nodes derive an identical secret without ever
transmitting it. It rests on one algebraic fact --- scalar multiplication commutes:
\begin{equation}
a \cdot (b \cdot G) \;=\; ab \cdot G \;=\; b \cdot (a \cdot G)
\end{equation}

\begin{figure}[H]
\centering
\begin{tikzpicture}[scale=1,every node/.style={transform shape}]
% headers
\node[node,minimum width=42mm,fill=mist] (u) at (-4.2,0) {UWS node U1};
\node[node,minimum width=42mm,fill=mist] (s) at ( 4.2,0) {SUB node S1};
% keys
\node[lbl,anchor=west,align=left] at (-6.4,-0.85)
  {private $a = \mathtt{7f3c\ldots91b4}$\ \ (secret)\\ public\ \ $A = a\cdot G$\ \ (curve point)};
\node[lbl,anchor=west,align=left] at ( 2.0,-0.85)
  {private $b = \mathtt{2d80\ldots4ce1}$\ \ (secret)\\ public\ \ $B = b\cdot G$\ \ (curve point)};
% open channel
\begin{scope}[on background layer]
\draw[draw=slate,dashed,rounded corners=2pt] (-6.6,-1.6) rectangle (6.6,-3.5);
\end{scope}
\node[font=\tiny\ttfamily,text=slate] at (0,-1.85)
  {OPEN ACOUSTIC CHANNEL --- THE ATTACKER SEES EVERYTHING IN THIS BAND};
\draw[flow] (-4.2,-2.35) -- node[lbl,above]{sends $A$} (4.2,-2.35);
\draw[flow] ( 4.2,-3.05) -- node[lbl,above]{sends $B$} (-4.2,-3.05);
% computation
\node[box,fill=mist,draw=teal,minimum width=40mm,minimum height=12mm] (cu) at (-4.2,-4.5)
  {$U1$ computes\\ $a\cdot B = ab\cdot G$};
\node[box,fill=mist,draw=teal,minimum width=40mm,minimum height=12mm] (cs) at ( 4.2,-4.5)
  {$S1$ computes\\ $b\cdot A = ba\cdot G$};
\node[box,fill=teal,draw=teal,text=white,minimum width=34mm,minimum height=12mm] (k) at (0,-4.5)
  {\textbf{SAME POINT $P$}\\ $K=\mathrm{SHA256}(P.x)$};
\draw[->,draw=teal,thick] (cu) -- (k);
\draw[->,draw=teal,thick] (cs) -- (k);
% attacker
\node[box,fill=palecoral,draw=coral,text=coral,minimum width=76mm] at (0,-5.9)
  {Attacker holds $A$ and $B$ but needs $a$ or $b$: this is the ECDLP, $\approx 2^{128}$ work};
\end{tikzpicture}
\caption[ECDH key agreement]{\textbf{ECDH key agreement.} Only public points cross the
channel. Each node multiplies the peer's point by its own private scalar and arrives at the
identical curve point $P$; hashing $P$'s $x$-coordinate yields a 256-bit symmetric key. The
key itself is never transmitted, so an eavesdropper who records the entire exchange still
has to solve the discrete logarithm problem.}
\label{fig:ecdh}
\end{figure}

In our simulator this is six lines of Python:

\begin{lstlisting}[style=py,caption={Key generation and ECDH shared-secret derivation
(\code{uwc\_simulation.py}, lines 109--118).}]
def gen_keys():
    pk = random.randint(1, curve.field.n - 1)   # private scalar
    return pk, pk * curve.g                     # public point

def shared(priv, pub):
    pt = priv * pub                             # ECDH: a.B on one side, b.A on the other
    return hashlib.sha256(str(pt.x).encode()).hexdigest()
\end{lstlisting}

\section{Nonces, timestamps and the replay attack}

A \textbf{nonce} is a ``number used once'': a fresh random value included in a message so
that no two runs of the protocol ever look alike. A \textbf{timestamp} records when the
message was created. Together they defeat the \textbf{replay attack}, which requires no
cryptography at all --- the adversary simply records a legitimate ciphertext and re-sends it
later. She cannot read it, but if the receiver accepts it she has impersonated a legitimate
node.

\begin{figure}[H]
\centering
\begin{tikzpicture}[x=1.06cm,y=1cm]
% acceptance band
\fill[palekelp] (0,0.1) rectangle (5,2.1);
\node[font=\scriptsize\bfseries,text=kelp] at (2.5,2.35) {ACCEPTANCE WINDOW $\Delta T = 5$\,s};
% axis
\draw[->,thick,draw=slate] (-0.6,0.1) -- (9.6,0.1) node[below left,font=\tiny\ttfamily]{time (s)};
\foreach \x/\lb in {0/0,1/1,2/2,3/3,4/4,5/5,6/6,7/7,8/8,9/9}
  {\draw[draw=slate] (\x,0.0) -- (\x,0.2); \node[font=\tiny\ttfamily,text=slate] at (\x,-0.28) {\lb};}
\draw[thick,draw=slate] (0,-0.05) -- (0,0.25);
\draw[thick,draw=coral,dashed] (5,0.1) -- (5,2.1);
\node[font=\tiny\ttfamily,text=coral,anchor=west] at (5.08,2.0) {$\Delta T$ boundary};
% legit send
\draw[thick,draw=slate] (0,0.1) -- (0,1.45);
\fill[slate] (0,1.45) circle (2.2pt);
\node[font=\tiny\ttfamily,anchor=west,align=left] at (0.15,1.72)
  {U1 sends $M_1$, $\TS=0.00$\\ (attacker records the ciphertext)};
% legit arrival
\draw[thick,draw=kelp] (1,0.1) -- (1,0.9);
\fill[kelp] (1,0.9) circle (2.2pt);
\node[font=\tiny\ttfamily,text=kelp,anchor=west] at (1.15,0.78)
  {arrives $t=1.0$: $|1.0-0.0|<5$ \ \textbf{ACCEPT}};
% replay
\draw[very thick,draw=coral,dashed] (6,0.1) -- (6,1.45);
\fill[coral] (6,1.45) circle (2.2pt);
\node[font=\tiny\ttfamily,text=coral,anchor=east,align=right] at (5.85,1.95)
  {attacker replays the same bytes at $t=6.0$;\\ the $\TS$ inside is still 0.00};
\node[box,fill=coral,draw=coral,text=white,font=\scriptsize\bfseries] at (7.6,1.05) {REJECT};
\draw[->,thick,draw=coral] (6.1,1.3) -- (7.15,1.15);
\node[font=\tiny\itshape,text=slate,anchor=west] at (-0.6,-0.85)
  {The timestamp sits \emph{inside} the encrypted payload, so the attacker cannot refresh it without the key.};
\end{tikzpicture}
\caption[Replay defeated by the freshness window]{\textbf{Why replay fails.} The receiver
compares its own clock against the timestamp carried inside the ciphertext. Because the
adversary cannot decrypt, she cannot update that timestamp; the stale value falls outside
$\Delta T$ and the message is dropped. The nonce is the second line of defence, catching a
replay that arrives while still inside the window. $\Delta T$ is measured in seconds here,
not milliseconds, precisely because acoustic propagation is slow.}
\label{fig:replay}
\end{figure}

\section{Mutual authentication}

One-way authentication means the server proves its identity --- that is what a browser
padlock provides. \textbf{Mutual} authentication means both sides prove identity to each
other. In a sensor network both directions matter: a buoy must not accept readings from a
counterfeit sensor, and a sensor must not surrender its data to a counterfeit buoy.

The proof always has the same shape: \emph{I send you a fresh challenge that only the real
you could answer, and you return it inside something only you could have produced.}

\section{The Dolev--Yao adversary}

Formal analysis assumes the strongest realistic adversary, named after Dolev and Yao
\cite{dolevyao}. She is not an attacker sitting on one link --- she \emph{is} the network.
She may read every message, delete messages, reorder them, replay old ones, and inject
anything she can construct from what she already knows. What she cannot do is break
cryptography: without the key she cannot decrypt, and without a private key she cannot
forge. This is the \emph{perfect cryptography} assumption.

\begin{figure}[H]
\centering
\begin{tikzpicture}
\node[node,minimum width=22mm] (a) at (-4.6,0) {Honest\\ node $A$};
\node[node,minimum width=22mm] (b) at ( 4.6,0) {Honest\\ node $B$};
\node[box,fill=palecoral,draw=coral,text=coral,minimum width=34mm,minimum height=13mm,
      font=\scriptsize\bfseries] (e) at (0,0) {ADVERSARY\\ \textnormal{\scriptsize is the channel}};
\draw[flow] (a) -- (e);
\draw[flow] (e) -- (b);
\node[box,anchor=north,align=left,minimum width=44mm,fill=palecoral,draw=coral] (can) at (-2.6,-1.5)
 {\textbf{\textcolor{coral}{CAN}}\\
  read every message\\ drop or delay any message\\ reorder messages\\
  replay old messages\\ inject forged messages};
\node[box,anchor=north,align=left,minimum width=44mm,fill=palekelp,draw=kelp] (cant) at (2.6,-1.5)
 {\textbf{\textcolor{kelp}{CANNOT}}\\
  decrypt without the key\\ forge without the private key\\
  invert a hash\\ guess a fresh nonce\\ break AES or the ECDLP};
\end{tikzpicture}
\caption[The Dolev--Yao adversary model]{\textbf{The Dolev--Yao adversary.} She controls the
entire network but not the mathematics. A protocol that survives this model has no flaw in
its \emph{message design}; any remaining weakness lies in the primitives or the
implementation. This is exactly the model Scyther checks against in
Chapter~\ref{ch:scyther}.}
\label{fig:dolevyao}
\end{figure}

\section{Notation used throughout}

\begin{table}[H]
\centering
\caption{Symbols and entities, following the base paper's Table~II.}
\label{tab:notation}
\small
\begin{tabularx}{\textwidth}{@{}L{26mm} X@{}}
\toprule
\textbf{Symbol} & \textbf{Meaning}\\
\midrule
UWS  & Underwater sensor: measures temperature, pressure, salinity, current velocity\\
SUB  & Submarine or AUV: mobile relay aggregating data from several sensors\\
B    & Surface buoy: bridges water to air, carries a GPS module\\
SAT  & Satellite: long-haul link from buoy to shore\\
BS   & Base station on land: final destination and trust anchor\\
\midrule
$\PK_i$    & Private key of entity $i$ (a secret scalar)\\
$\PU_i$    & Public key of entity $i$, $\PU_i = \PK_i \times \rho$ (a curve point)\\
$\rho$     & Generator point of the curve\\
$\ID_i$    & Identity of $i$, derived as a hash of $\PU_i$\\
$\RID_i$   & Registration identifier, $H(\ID_i \cat \PU_i \cat \PK_i)$\\
$\vartheta_i$ & Verification value, $H(\RID_i \cat \PU_i)$\\
$\Nonce_i$ & Fresh random number generated by $i$ for one run\\
$\TS_n$    & Timestamp attached to message $n$\\
$\Delta T$ & Freshness window; accept only if $|\mathrm{now} - \TS| < \Delta T$\\
$E\{\PU_j, \langle m\rangle\}$ & $m$ encrypted so only the holder of $\PK_j$ can read it\\
$D\{\PK_j, (\cdot)\}$ & Decryption with the private key of $j$\\
$\cat$     & Concatenation of fields\\
$R$        & Set of registered $\RID$s held by a node (hash table or Bloom filter)\\
$S_{i\leftrightarrow j}$ & Session established between $i$ and $j$\\
ENV        & Aggregated environmental payload $\{$Temp, Press, CurrVel, Sal$\}$\\
\bottomrule
\end{tabularx}
\end{table}
